\documentclass[12pt]{article}
\usepackage[utf8]{inputenc}
\usepackage[margin=1in]{geometry}
\usepackage{setspace}
\usepackage{hyperref}  % Allows for clickable links
\setstretch{1.2}

\begin{document}

\title{The Cost of Liberation}
\author{ }
\date{ }
\maketitle

\noindent
\textbf{The Cost of Liberation}

\medskip

The idea of true freedom within these United States is steeped in a profound irony. Although the freed citizen sees it as a triumph of freeing the country, and ourselves, from the rule of Great Britain, within the short 70 year time period from the Declaration of Independence, the country had begun to indulge in the most anti-freedom measures imaginable, arguably even worse than the measures Great Britain had inflicted upon the country all those years ago. For those who remained unfree, the cost for their freedom would be just as steep as it was for the founders of the country that now imprisoned them: life itself. However, for many, the cost would not be mere existence, but also include a litany of other costs.\cite{Foner1999, Berlin1998}

\medskip

For Northerners, the Civil War primarily centered on preserving the Union rather than emancipation. Many moderates believed conflict should have been avoided altogether. They paid a psychological and moral price upon realizing that the war had evolved into a confrontation over slavery.\cite{McPherson1988} The Civil War was one of the bloodiest battles in US history, with many soldiers returning to their families a scarred person.\cite{Faust2009} When the Emancipation Proclamation was issued in 1863, some Union soldiers questioned having signed up just to free enslaved people.\cite{Lincoln1863,McPhersonFreed}

\medskip

In the Confederacy, the war was seen as a way to maintain their “way of life.” Plantation owners’ wealth hinged on enslavement, so the end of slavery meant an imminent financial crisis for them.\cite{Phillips1929,Wright1978} Beyond the slave owners (only sub-50\% of families owned slaves), many Southerners strongly believed that black people were inherently inferior, almost less than human, and they were seen as savages (rebellions such as Turner’s propagated this belief).\cite{1860Census,Stampp1956} Therefore, losing the war meant economic ruin and an inner crisis, as the entire bedrock for their worldview was being dismantled. In their eyes, freedom for the enslaved effectively meant the end of their lives as they knew it. Perhaps they were more right than they knew.

\medskip

While soldiers fought and died in the thousands, enslaved people had long been paying in blood through forced labor, physical abuse, and the omnipresent threat of being sold away from their families. Enslaved families could be torn apart at any moment, sold to different owners. Harriet Jacobs, in \textit{Incidents in the Life of a Slave Girl}, endured psychological torment in a tiny attic for seven years to protect her children, highlighting the terror and sacrifices enslaved mothers faced.\cite{Jacobs1861,Yellin2004} Many enslaved people lived in constant cognitive dissonance, forced to attend church services where pastors twisted scripture to uphold slavery. Frederick Douglass famously denounced the hypocrisy of the “Christianity of this land,” contrasting it with the true teachings of Christ.\cite{Douglass1845, DouglassChurch}

\medskip

Even after the Emancipation Proclamation, Black Americans quickly learned that ‘technical’ liberty did not translate to acceptance or equality in the eyes of most white Americans.\cite{Franklin2000, Blight2018} Despite newfound freedom, newly liberated people found themselves isolated in a society that refused to treat them as equals. Sharecropping replaced slavery, keeping formerly enslaved families locked in cycles of debt, and Jim Crow laws emerged to enforce racial segregation, ensuring that Black Americans would face discrimination in every facet of life.\cite{Foner1988, Woodward1955}

\medskip

The Union sacrificed thousands of soldiers, driven by an evolving moral understanding that slavery was incompatible with national ideals. The Confederacy saw its social and economic order crumble, paying a price in lives, finances, and moral self-image.\cite{Escott1978, Rable1991} Even when emancipation came, at the cost of a bloody civil war and many lives, the freed slaves discovered that freedom did not guarantee equality. Instead, they entered another century of systemic discrimination and segregation, only seeing tangible hope of legal and social acceptance during the Civil Rights Movement in the ‘60s.\cite{Branch1989, Carson1981} Less than 70 years ago, the first black students graduated from USC, truly showcasing this painstaking, uneven progress even on this campus.\cite{USCTimeline}

\vfill

\section*{References}
\begin{thebibliography}{99}

\bibitem{Foner1999}
Foner, Eric. 
\textit{The Story of American Freedom}. 
New York: W. W. Norton \& Company, 1999.  
\href{https://wwnorton.com/books/9780393319620}{https://wwnorton.com/books/9780393319620}

\bibitem{Berlin1998}
Berlin, Ira. 
\textit{Many Thousands Gone: The First Two Centuries of Slavery in North America}. 
Cambridge, MA: Belknap Press, 1998.  
\href{https://www.hup.harvard.edu/catalog.php?isbn=9780674002111}{https://www.hup.harvard.edu/catalog.php?isbn=9780674002111}

\bibitem{McPherson1988}
McPherson, James M. 
\textit{Battle Cry of Freedom: The Civil War Era}. 
New York: Oxford University Press, 1988.  
\href{https://global.oup.com/academic/product/battle-cry-of-freedom-9780195038637}{https://global.oup.com/academic/product/battle-cry-of-freedom-9780195038637}

\bibitem{Faust2009}
Faust, Drew Gilpin. 
\textit{This Republic of Suffering: Death and the American Civil War}.
New York: Vintage, 2009.  
\href{https://www.penguinrandomhouse.com/books/50278/this-republic-of-suffering-by-drew-gilpin-faust/}{https://www.penguinrandomhouse.com/books/50278/this-republic-of-suffering-by-drew-gilpin-faust/}

\bibitem{Lincoln1863}
Lincoln, Abraham. 
``The Emancipation Proclamation,'' January 1, 1863.  
\href{https://www.archives.gov/exhibits/featured-documents/emancipation-proclamation}{https://www.archives.gov/exhibits/featured-documents/emancipation-proclamation}

\bibitem{McPhersonFreed}
McPherson, James M. 
``Who Freed the Slaves?'' 
\textit{Proceedings of the American Philosophical Society} 139, no. 1 (1995): 1–10.  
\href{https://www.jstor.org/stable/987110}{https://www.jstor.org/stable/987110}

\bibitem{Phillips1929}
Phillips, Ulrich Bonnell. 
\textit{Life and Labor in the Old South}. 
Boston: Little, Brown and Company, 1929.  
\href{https://archive.org/details/lifelaborinolds00phil}{https://archive.org/details/lifelaborinolds00phil}

\bibitem{Wright1978}
Wright, Gavin. 
\textit{The Political Economy of the Cotton South}. 
New York: W. W. Norton \& Co., 1978.  
\href{https://wwnorton.com/books/9780393090918}{https://wwnorton.com/books/9780393090918}

\bibitem{1860Census}
U.S. Census Bureau. 
``1860 Census.''  
\href{https://www.census.gov/history/www/through_the_decades/overview/1860.html}{https://www.census.gov/history/www/through_the_decades/overview/1860.html}

\bibitem{Stampp1956}
Stampp, Kenneth M. 
\textit{The Peculiar Institution: Slavery in the Ante-Bellum South}. 
New York: Vintage, 1956.  
\href{https://www.penguinrandomhouse.com/books/170303/the-peculiar-institution-by-kenneth-m-stampp/}{https://www.penguinrandomhouse.com/books/170303/the-peculiar-institution-by-kenneth-m-stampp/}

\bibitem{Jacobs1861}
Jacobs, Harriet. 
\textit{Incidents in the Life of a Slave Girl}. 
Boston: Published for the Author, 1861.  
\href{https://www.gutenberg.org/ebooks/11030}{https://www.gutenberg.org/ebooks/11030}

\bibitem{Yellin2004}
Yellin, Jean Fagan. 
\textit{Harriet Jacobs: A Life}. 
New York: Basic Civitas Books, 2004.  
\href{https://www.basicbooks.com/titles/jean-fagan-yellin/harriet-jacobs/9780465092887/}{https://www.basicbooks.com/titles/jean-fagan-yellin/harriet-jacobs/9780465092887/}

\bibitem{Douglass1845}
Douglass, Frederick. 
\textit{Narrative of the Life of Frederick Douglass, an American Slave, Written by Himself}.
Boston: Anti-Slavery Office, 1845.  
\href{https://www.gutenberg.org/ebooks/23}{https://www.gutenberg.org/ebooks/23}

\bibitem{DouglassChurch}
Douglass, Frederick. 
``The Church and Prejudice.'' 
(Various speeches denouncing religious justifications of slavery, 1841--1860).  
\href{https://rbscp.lib.rochester.edu/2945}{https://rbscp.lib.rochester.edu/2945}

\bibitem{Franklin2000}
Franklin, John Hope, and Loren Schweninger.
\textit{Runaway Slaves: Rebels on the Plantation}.
New York: Oxford University Press, 2000.  
\href{https://global.oup.com/academic/product/runaway-slaves-9780195084511}{https://global.oup.com/academic/product/runaway-slaves-9780195084511}

\bibitem{Blight2018}
Blight, David W.
\textit{Frederick Douglass: Prophet of Freedom}.
New York: Simon \& Schuster, 2018.  
\href{https://www.simonandschuster.com/books/Frederick-Douglass/David-W-Blight/9781416590323}{https://www.simonandschuster.com/books/Frederick-Douglass/David-W-Blight/9781416590323}

\bibitem{Foner1988}
Foner, Eric.
\textit{Reconstruction: America's Unfinished Revolution, 1863--1877}.
New York: Harper \& Row, 1988.  
\href{https://www.harpercollins.com/products/reconstruction-updated-edition-eric-foner}{https://www.harpercollins.com/products/reconstruction-updated-edition-eric-foner}

\bibitem{Woodward1955}
Woodward, C. Vann.
\textit{The Strange Career of Jim Crow}.
New York: Oxford University Press, 1955.  
\href{https://global.oup.com/academic/product/the-strange-career-of-jim-crow-9780195146905}{https://global.oup.com/academic/product/the-strange-career-of-jim-crow-9780195146905}

\bibitem{Escott1978}
Escott, Paul D.
\textit{After Secession: Jefferson Davis and the Failure of Confederate Nationalism}.
Baton Rouge: LSU Press, 1978.  
\href{https://lsupress.org/books/detail/after-secession/}{https://lsupress.org/books/detail/after-secession/}

\bibitem{Rable1991}
Rable, George C.
\textit{Civil Wars: Women and the Crisis of Southern Nationalism}.
Urbana: University of Illinois Press, 1991.  
\href{https://www.press.uillinois.edu/books/catalog/84qmf3gp9780252062060.html}{https://www.press.uillinois.edu/books/catalog/84qmf3gp9780252062060.html}

\bibitem{Branch1989}
Branch, Taylor.
\textit{Parting the Waters: America in the King Years 1954--63}.
New York: Simon \& Schuster, 1989.  
\href{https://www.simonandschuster.com/books/Parting-the-Waters/Taylor-Branch/9780671687427}{https://www.simonandschuster.com/books/Parting-the-Waters/Taylor-Branch/9780671687427}

\bibitem{Carson1981}
Carson, Clayborne.
\textit{In Struggle: SNCC and the Black Awakening of the 1960s}.
Cambridge, MA: Harvard University Press, 1981.  
\href{https://www.hup.harvard.edu/catalog.php?isbn=9780674447271}{https://www.hup.harvard.edu/catalog.php?isbn=9780674447271}

\bibitem{USCTimeline}
University of South Carolina Archives. 

\end{thebibliography}

\end{document}
